\documentclass[letterpaper,12pt]{article}

\usepackage[printonlyused]{acronym}
\usepackage{amsmath}
\usepackage[font={small,sf},labelfont={small,sf}]{caption}
\usepackage{color}
\usepackage{csquotes}
\usepackage{fancyhdr}
\usepackage[includeheadfoot,left=1in,top=.4in,right=1in,bottom=.75in,headsep=\dimexpr3cm-59pt\relax,headheight=59pt]{geometry}
\usepackage{graphicx}
\usepackage[pagebackref,hyperindex=true]{hyperref}
\usepackage{listings}
\usepackage{longtable}
\usepackage{mdwlist}
\usepackage{parskip}
\usepackage{setspace}
\usepackage{tabularx}
\usepackage[compact]{titlesec}
\usepackage{xfrac}

% default font packages
\usepackage{courier} % use courier for the mono-spaced font 
\usepackage{helvet}  % use a Helvetica clone for default text (sans-serif)

%%%
% default document settings
%%%

% Proposal metadata has its own namespace so it cannot replace LaTeX's
% standard metadata setters.  Configure these values with \renewcommand after
% loading this template.  \newcommand intentionally reports an error if this
% template (or another package) has already declared one of these names.
\newcommand{\ProposalProjectName}{Untitled Project}
\newcommand{\ProposalTitle}{\ProposalProjectName}
\newcommand{\ProposalAuthor}{}
\newcommand{\ProposalCompany}{}

% Keep the standard LaTeX metadata in sync for packages and commands which use
% \@title and \@author (for example, PDF metadata and \maketitle).
\title{\ProposalTitle}
\author{\ProposalAuthor}

% setup the default fonts for section headers
\titleformat*{\section}{\sffamily\fontsize{16}{16}\selectfont\bfseries}
\titleformat*{\subsection}{\sffamily\fontsize{14}{14}\selectfont\bfseries}
\titleformat*{\subsubsection}{\sffamily\fontsize{12}{12}\selectfont\bfseries}
\titleformat*{\paragraph}{\sffamily\fontsize{12}{12}\selectfont\bfseries}
\titleformat*{\subparagraph}{\sffamily\fontsize{12}{12}\selectfont\bfseries}

% section numbering
\setcounter{tocdepth}{3}     % display only 3 sections deep in the table of contents
\setcounter{secnumdepth}{5}  % number to 5 sections deep

% add acronyms to the TOC (use chapter, if chapters are available, otherwise use sections
% Based off of suggestions at: http://jevopi.blogspot.com/2009/09/acronyms-and-latex.html
\providecommand{\listofacronymsname}{List of Acronyms}
\providecommand{\listofacronyms}{
    \ifcsname chapter\endcsname
        \chapter*{\listofacronymsname}
        \addcontentsline{toc}{chapter}{\listofacronymsname}
    \else
        \section*{\listofacronymsname}
        \addcontentsline{toc}{section}{\listofacronymsname}
    \fi
    \label{sec:acronyms}
	\markboth{\listofacronymsname}{\listofacronymsname}
    \input{includes/acronyms}
}

%%%
% custom page styles
%%%

\fancypagestyle{restricted}{
    \fancyhead[L]{\rule[-2ex]{0pt}{2ex}\small \MakeUppercase{\baa}}
    \fancyhead[c]{ 
        \small
        \ifdefined\classification
        \textbf{\classification} \\ 
        \fi
        \textbf{ \ProposalCompany } \\
        \textbf{\MakeUppercase{\restrictions} } \\
        \ProposalTitle: \biline
    }
    \fancyhead[R]{\small \MakeUppercase{\baa}}

    \fancyfoot[L]{\small \ProposalTitle}
    \fancyfoot[c]{\small \textbf{{\MakeUppercase{\restrictions}}}}
    \fancyfoot[c]{
        \small
        \textbf{\MakeUppercase{\restrictions} } \\
        \ifdefined\classification
        \textbf{\classification}
        \fi
    }
    \fancyfoot[R]{\small \thepage}
}

\fancypagestyle{plain}{
    \fancyhead[L]{\rule[-2ex]{0pt}{2ex}\small \MakeUppercase{\baa}}
    \fancyhead[C]{ 
        \small
        \ifdefined\classification
        \textbf{\classification} \\ 
        \fi
        \textbf{ \ProposalCompany } \\
        \ProposalTitle: \biline
    }
    \fancyhead[R]{\small \MakeUppercase{\baa}}

    \fancyfoot[L]{\small \ProposalTitle}
    \fancyfoot[C]{
        \small
        \ifdefined\classification
        \textbf{\classification}
        \fi
    }
    \fancyfoot[R]{\small \thepage}
}

% set the default page style to "plain"
\pagestyle{plain}

%%%
% page templates
%%%

\newcommand{\whitepapercover}{
    \begin{titlepage}
    \thispagestyle{empty}
    \vspace*{1.25in}

    \hfill \textsc{\sffamily\huge\bfseries \ProposalTitle} \\
    \begin{flushright}
        \textsc{\sffamily\large \biline} \\
        \textsc{\sffamily\large \today}
    \end{flushright}

    \vspace{3.0in}
    \begin{center}
        \includegraphics[width=6.0in]{includes/logo.png}
    \end{center}

    \end{titlepage}
    \cleardoublepage
}

\newcommand{\proposalcover}{
    \thispagestyle{empty}
    \vspace*{-1.25in}
    \begin{center}
        \includegraphics[width=4.0in]{includes/logo.png}
    \end{center}
  
    {
        \footnotesize
        \begin{center}
	        \begin{tabular}{|l|p{3.5in}|}
	            \hline
	            \textbf{BAA Number} & \baa \\ \hline
	            \textbf{Technical Area} & \techarea \\ \hline
	            \textbf{Lead Organization} & \ProposalCompany \\ \hline
                
                \ifdefined\companytype
	            \textbf{Type of Business} & \companytype \\ \hline
                \fi

                \ifdefined\companyref
	            \textbf{Contractor's Reference Number} & \companyref \\ \hline
                \fi

                \ifdefined\team
	            \textbf{Team Members} & \team \\ \hline
                \fi

	            \textbf{Proposal Title} & \ProposalTitle: \biline \\ \hline
	            \textbf{Technical Point of Contact} & \ProposalAuthor \\ & \email \\ & \phone \\ & \address \\ \hline
	            \textbf{Administrative Point of Contact} & \ProposalAuthor \\ & \email \\ & \phone \\ & \address \\ \hline

                \ifdefined\cost
	            \textbf{Proposed Cost} & \cost \\ \hline
                \fi

                \ifdefined\submitdate
	            \textbf{Date Submitted} & \submitdate \\ \hline
                \fi

                \ifdefined\awardtype
	            \textbf{Award Instrument Requested} & \awardtype \\ \hline
                \fi

                \ifdefined\placeofperformance
	            \textbf{Place of Performance} & \placeofperformance \\ \hline
                \fi

                \ifdefined\pop
	            \textbf{Period of Performance}  & \pop \\ \hline
                \fi

	            \textbf{Date Prepared} & \today \\ \hline

                \ifdefined\duns
	            \textbf{DUNS Number} & \duns \\ \hline
                \fi

                \ifdefined\tin
	            \textbf{TIN number} & \tin \\ \hline
                \fi
                
                \ifdefined\cagecode
	            \textbf{Cage Code} & \cagecode \\ \hline
                \fi
	            \textbf{Proposal Validity Period} & This proposal is valid for 120 days \\ \hline
	        \end{tabular}
        \end{center}
    }

    \ifproposalrestricted
    \section*{Use and Disclosure of Data}
    \proposalrestrictedlegend
    \fi
    \cleardoublepage
}

\newcommand{\bib}{
    \cleardoublepage
    \pagestyle{plain}
    \phantomsection
    \ifcsname chapter\endcsname
        \addcontentsline{toc}{chapter}{References}
    \else
        \addcontentsline{toc}{section}{References}
    \fi
    \bibliographystyle{unsrt} % unsrt = plain, except sorted by use, not date.
    {
        \raggedright
        \bibliography{includes/refs}
    }
}

\newcommand{\docinfo}{
    \pagenumbering{roman}
    \thispagestyle{plain}
    \section*{Copyright}
    Copyright \copyright \the\year\ \ProposalCompany.

    All trademarks within this document belong to their respective owners.

    \section*{Authors}
    \ProposalAuthor, \ProposalCompany

    \section*{Point of Contact}
            
    \ProposalAuthor \\ \email \\ \phone
    
    \ifexportcontrolled
    \section*{Export Control}
    \exportcontrollegend
    \fi
    
    \cleardoublepage

    \tableofcontents
    \listoftables
    \listoffigures
    \listofacronyms 
    \cleardoublepage
    \pagenumbering{arabic}
}

%%%
% ease of use macros
%%%

% Example: \q{foo}
% Results: Language-sensitive quotation marks around "foo".
\newcommand{\q}[1]{\enquote{#1}}

% example: C:$\bs$Program Files$\bs$Adobe$\bs$
% Results: C:\Program Files\Adobe\
\def \bs{\char`\\}

% example: \begin{fig}{figure label}{figure caption}{ ... }
% results: A figure, boxed in the center, with font slightly shrunk, with a label and caption
\newcommand{\temporarylabel}{}
\newcommand{\temporarycaption}{}
\newenvironment{fig}[2]{
    \renewcommand{\temporarylabel}{#1}
    \renewcommand{\temporarycaption}{#2}
    \begin{figure}[!htbp]
    \begin{center}
    \begin{small}
}{
    \end{small}
    \end{center}
    \caption{\temporarycaption \label{\temporarylabel}}
    \end{figure}
}
